% ===========================================================================
\section{Two channels where noise is free}\label{sec:first-examples}
\warmth{0}\quad\whisper{Noise is harmless as long as it never makes two inputs look alike.}

\begin{strand}
Start with the two easy extremes. If the channel is noiseless, one error-free
bit per use is obvious. If the channel is noisy but different inputs produce
disjoint sets of outputs, the receiver still recovers the input exactly --- so
the capacity is again that of a clean pipe. What costs capacity is not
randomness, but \emph{confusability}.
\end{strand}

\trigger{Before computing anything, ask whether the output ever leaves the input in doubt.}

\block{The noiseless binary channel}

\begin{example}[noiseless binary channel]\label{ex:noiseless-binary}
Let $\X=\Y=\{0,1\}$ and let every input be reproduced exactly.
\begin{center}
\begin{tikzpicture}[x=2.4cm,y=1.1cm,>=latex,font=\small]
  \node (x0) at (0,1) {$0$};
  \node (x1) at (0,0) {$1$};
  \node (y0) at (1,1) {$0$};
  \node (y1) at (1,0) {$1$};
  \draw[->,indigo] (x0) -- (y0);
  \draw[->,indigo] (x1) -- (y1);
\end{tikzpicture}
\end{center}
One error-free bit is transmitted per use of the channel, so we expect $C=1$.
The definition confirms it: here $Y=X$, hence $\Ent(Y\mid X)=0$ and
\[
  C=\max_{p(x)}\MI(X;Y)
   =\max_{p(x)}\Ent(X)
   =\answerin[1.2cm]{1}\ \text{bit},
\]
the maximum being attained at $p(x)=\answerin[2.4cm]{(\tfrac12,\tfrac12)}$.
\studentrecall{$C=1$ bit, attained by the uniform input $(\tfrac12,\tfrac12)$.}
\end{example}

\block{A noisy channel with non-overlapping outputs}

\begin{example}[non-overlapping outputs]\label{ex:nonoverlapping}
Let $\X=\{0,1\}$ and $\Y=\{1,2,3,4\}$, with
\[
  p(1\mid0)=p(2\mid0)=\tfrac12,
  \qquad
  p(3\mid1)=\tfrac13,\quad p(4\mid1)=\tfrac23 .
\]
\begin{center}
\begin{tikzpicture}[x=2.6cm,y=0.95cm,>=latex,font=\small]
  \node (x0) at (0,2.6) {$0$};
  \node (x1) at (0,0.4) {$1$};
  \node (y1) at (1,3.3) {$1$};
  \node (y2) at (1,2.1) {$2$};
  \node (y3) at (1,0.9) {$3$};
  \node (y4) at (1,-0.3) {$4$};
  \draw[->,indigo] (x0) -- node[above,sloped,font=\scriptsize]{$1/2$} (y1);
  \draw[->,indigo] (x0) -- node[below,sloped,font=\scriptsize]{$1/2$} (y2);
  \draw[->,madder] (x1) -- node[above,sloped,font=\scriptsize]{$1/3$} (y3);
  \draw[->,madder] (x1) -- node[below,sloped,font=\scriptsize]{$2/3$} (y4);
\end{tikzpicture}
\end{center}
The channel is genuinely noisy: the output is random for both inputs. But the
output sets $\{1,2\}$ and $\{3,4\}$ are \keyword{disjoint}, so from $y$ the
receiver can always tell which input was sent. Hence
\[
  \Ent(X\mid Y)=\answerin[1.2cm]{0},
  \qquad
  \MI(X;Y)=\Ent(X)-\Ent(X\mid Y)=\Ent(X),
\]
\studentrecall{Disjoint output sets mean $Y$ determines $X$, so $\Ent(X\mid Y)=0$ and $\MI=\Ent(X)$.}
and therefore
\[
  C=\max_{p(x)}\Ent(X)=1\ \text{bit},
\]
again attained by the uniform input. Randomness inside each row is irrelevant;
only overlap between rows costs capacity.
\end{example}

\begin{remark}
It is worth naming the general principle: the capacity of a channel depends on
its transition matrix only through which inputs can be \emph{confused}. If the
supports of distinct rows are pairwise disjoint, the channel is as good as a
noiseless channel on $|\X|$ symbols and $C=\log|\X|$.
\end{remark}

\begin{yourturn}
\keyword{Noisy typewriter.} Let $\X=\Y=\{A,B,\ldots,Z\}$ (26 letters) and
suppose each letter is either printed correctly or printed as the next letter,
each with probability $1/2$. By using only every other letter of the alphabet as
an input, the outputs of the used inputs become disjoint. What capacity does
this suggest, and why can it not be beaten?
\studentrecall{$13$ usable non-confusable inputs give $C=\log13=\log26-1$ bits.}
\ifshowanswers\else\workspace[3]\fi
\answerprose{Using $13$ letters spaced two apart makes the output sets disjoint,
so those $13$ inputs form a noiseless channel and $C\ge\log13$. Conversely each
row has entropy $\Ent(Y\mid X=x)=1$ bit, so
$\MI(X;Y)=\Ent(Y)-1\le\log26-1=\log13$. Hence $C=\log13$ bits, attained by the
uniform distribution on all $26$ inputs as well.}
\end{yourturn}

\recall{Which feature of the transition matrix actually destroys capacity: randomness within rows, or overlap between rows?}
